\begin{frame}{마르코프 성질 (Markov property)}
  \begin{block}{정의}
    \kb{마르코프 성질}: 다음 상태는 오직 \textbf{현재} 상태와 행동에만
    의존한다 --- 어떻게 지금 상태에 도달했는지(과거 전체 이력)는
    상관없다:
    \[
    P(s_{t+1} \mid s_t, a_t, s_{t-1}, a_{t-1}, \ldots, s_0)
    \;=\; P(s_{t+1} \mid s_t, a_t)
    \]
  \end{block}
  \vskip0.6em
  \begin{itemize}
    \item \textbf{예}: 체스판의 현재 배치만 알면, 거기까지 어떤
      수순을 거쳐왔는지는 다음 수를 정하는 데 필요 없다.
    \item 이 성질 덕분에 ``지금까지의 전체 이력''이 아니라
      ``현재 상태'' \textbf{하나만 기억하면 충분}해진다 --- 뒤에서
      배울 알고리즘들이 상태 하나만 보고 결정을 내릴 수 있는
      이유가 여기서 나온다.
  \end{itemize}
  \vskip0.5em
  \textbf{``상관없다''가 확률적으로 정확히 무엇을 뜻하는가} ---
  조건부 확률의 정의와 사건의 독립성으로부터:
  \[
  P(s_{t+1}=s' \mid s_t, a_t, s_{t-1}, a_{t-1}, \ldots)
  \;=\; P(s_{t+1}=s' \mid s_t, a_t)
  \]
  가 \textbf{모든} 과거 이력과 \textbf{모든} 다음 상태 $s'$ 에 대해
  성립할 때, $s_t$ 가 ($a_t$ 를 준 상태에서) \kb{마르코프 성질}을
  가른다고 한다.
  \vskip0.4em
  {\footnotesize ``상관없다''는 직관이 이 식으로 번역되는 순간,
  \textbf{검증 가능한 주장}이 된다 --- 실습 노트북에서 실제로
  샘플링해서 양변의 확률추정치가 일치하는지 확인한다.}
\end{frame}
